\documentclass[11pt]{article}

\usepackage[T1]{fontenc}
\usepackage[utf8]{inputenc}
\usepackage[margin=1in]{geometry}
\usepackage{amsmath,amssymb,amsthm,mathtools}
\usepackage{enumitem}
\usepackage{booktabs}
\usepackage{longtable}
\usepackage{array}
\usepackage[hidelinks]{hyperref}

\newtheorem{definition}{Definition}[section]
\newtheorem{lemma}[definition]{Lemma}
\newtheorem{proposition}[definition]{Proposition}
\newtheorem{theorem}[definition]{Theorem}
\newtheorem{remark}[definition]{Remark}
\newtheorem{convention}[definition]{Convention}

\newcommand{\R}{R}
\newcommand{\Sfam}{\mathcal S}
\newcommand{\pow}[1]{R^{#1}}
\newcommand{\Graph}{\operatorname{Graph}}
\newcommand{\Imf}{\operatorname{Im}}
\newcommand{\pr}{\operatorname{pr}}
\newcommand{\dom}{\operatorname{dom}}
\newcommand{\LC}{\operatorname{LC}}
\newcommand{\Cont}{\operatorname{Cont}}
\newcommand{\LInc}{\operatorname{LInc}}
\newcommand{\LDec}{\operatorname{LDec}}
\newcommand{\val}{\operatorname{val}}
\newcommand{\nhd}{\operatorname{nhd}}
\newcommand{\core}{\operatorname{core}}
\newcommand{\cat}{\mathbin{{}^\frown}}
\newcommand{\Lean}[1]{\texttt{#1}}

\title{Detailed mathematical write-up of the Lean formalisation\\
O-minimal structures, definable maps, images, and local properties}
\author{}
\date{}

\begin{document}
\maketitle

\section{Purpose and high-level structure}

This document explains the Lean development in mathematical language.  The Lean file builds an
axiomatic version of an o-minimal structure over a densely ordered set, defines definable sets and
definable functions, proves that several set-theoretic constructions preserve definability, and then
uses these constructions to show that certain loci of one-variable definable functions are definable.

The additional result added to the Lean file is the following standard theorem.

\begin{theorem}[Image of a definable function]
Let $M$ be an o-minimal structure on a densely ordered set $R$.  Let
\[
  A\subseteq R^m,
  \qquad
  B\subseteq R^n,
\]
and let
\[
  f:A\to B
\]
be a definable function.  Then
\[
  \Imf(f)=\{y\in R^n: \exists x\in A,\; y=f(x)\}
\]
is a definable subset of $R^n$.  Moreover,
\[
  \Imf(f)\subseteq B.
\]
\end{theorem}

The proof of this theorem does not use the special one-dimensional o-minimality axiom.  It only uses
the closure of definable sets under coordinate reindexing and existential projection.  This matches
the Lean proof: the graph of $f$ is reindexed from coordinates $(x,y)$ to coordinates $(y,x)$ and then
projected onto the $y$-coordinates.

\section{Tuples, coordinates, and basic set operations}

\subsection{Tuple spaces}

For each natural number $n$, the Lean file writes
\[
  \Lean{Power R n}
\]
for the type of $n$-tuples of elements of $R$.  Formally this is implemented as
\[
  \operatorname{Fin}(n)\to R.
\]
This representation is convenient in Lean because coordinates are indexed by the finite type
$\operatorname{Fin}(n)$.

Mathematically, we write
\[
  \Lean{Power R n}=R^n.
\]
If $x\in R^n$, then $x_i$ denotes the $i$-th coordinate of $x$.

\subsection{Concatenation and block projections}

If
\[
  x\in R^n,
  \qquad
  y\in R^m,
\]
then their concatenation is written
\[
  x\cat y\in R^{n+m}.
\]
In Lean this is
\[
  \Lean{Power.append x y}.
\]

If $z\in R^{n+m}$, then the left and right blocks are denoted by
\[
  z_L\in R^n,
  \qquad
  z_R\in R^m.
\]
In Lean these are
\[
  \Lean{Power.left z},
  \qquad
  \Lean{Power.right z}.
\]
They satisfy the expected identities
\[
  (x\cat y)_L=x,
  \qquad
  (x\cat y)_R=y.
\]
These are the Lean simp lemmas
\[
  \Lean{Power.left\_append},
  \qquad
  \Lean{Power.right\_append}.
\]

\subsection{Basic Boolean operations on sets}

A set in Lean is represented as a predicate.  Thus a subset $A\subseteq R^n$ is implemented as
\[
  \Lean{Set (Power R n)}.
\]
The file defines the following elementary set operations directly as predicates:
\[
\begin{aligned}
  \Lean{setEmpty} &= \emptyset,\\
  \Lean{setUnion A B} &= A\cup B,\\
  \Lean{setInter A B} &= A\cap B,\\
  \Lean{setCompl A} &= (R^n\setminus A).
\end{aligned}
\]
For example,
\[
  (\Lean{setInter A B})(x)
  \quad\Longleftrightarrow\quad
  A(x)\land B(x).
\]

\section{The ordered base set}

\begin{definition}[Dense linear order without endpoints]
A dense linear order without endpoints consists of a type $R$ together with a binary relation $<$ such
that the following properties hold.
\begin{enumerate}[label=(\roman*)]
  \item Irreflexivity: $\neg(x<x)$ for every $x\in R$.
  \item Transitivity: if $x<y$ and $y<z$, then $x<z$.
  \item Trichotomy: for all $x,y\in R$, exactly one of $x<y$, $x=y$, or $y<x$ holds.
  \item Density: if $x<y$, then there exists $z\in R$ such that $x<z<y$.
  \item No left endpoint: for every $x\in R$, there exists $y\in R$ such that $y<x$.
  \item No right endpoint: for every $x\in R$, there exists $z\in R$ such that $x<z$.
\end{enumerate}
In Lean this is the structure
\[
  \Lean{DenseLinearOrderNoEndpoints}.
\]
\end{definition}

The Lean development does not use the full power of an imported order typeclass.  Instead, the order
relation and the required axioms are packaged manually in the structure above.  This keeps the
formalisation self-contained.

\section{Intervals and the one-dimensional o-minimality axiom}

\subsection{Endpoints}

To state intervals uniformly, the Lean file introduces an endpoint type
\[
  R\cup\{-\infty,+\infty\}.
\]
In Lean this is the inductive type
\[
  \Lean{Endpoint R},
\]
with constructors
\[
  \Lean{Endpoint.negInf},
  \qquad
  \Lean{Endpoint.finite a},
  \qquad
  \Lean{Endpoint.posInf}.
\]
The endpoint order is the obvious extension of the order on $R$:
\[
  -\infty<a<+\infty
\]
for every finite $a\in R$.

\subsection{Points and intervals}

For $a\in R$, the singleton $\{a\}\subseteq R$ is represented as a subset of $R^1$:
\[
  \Lean{pointSet a}=\bigl\{x\in R^1:x_0=a\bigr\}.
\]

For endpoints $a,b\in R\cup\{-\infty,+\infty\}$, the open interval is
\[
  (a,b)=\{x\in R:a<x<b\},
\]
with the usual interpretation when one endpoint is infinite.  In Lean this is
\[
  \Lean{openInterval D a b}.
\]

\begin{definition}[Finite unions of points and intervals]
A subset $A\subseteq R$ is a finite union of points and intervals if it belongs to the smallest class
containing:
\begin{enumerate}[label=(\roman*)]
  \item the empty set;
  \item every singleton $\{a\}$;
  \item every open interval $(a,b)$, with possibly infinite endpoints;
  \item the union of two members of the class.
\end{enumerate}
In Lean this is the inductive predicate
\[
  \Lean{FiniteUnionOfPointsAndIntervals D A}.
\]
\end{definition}

\section{O-minimal structures in the Lean file}

\begin{definition}[Axiomatic o-minimal structure]
Let $(R,<)$ be a dense linear order without endpoints.  An o-minimal structure on $R$, as formalised
in the Lean file, consists of a family
\[
  \Sfam_n\subseteq \mathcal P(R^n),
  \qquad n=0,1,2,\ldots,
\]
satisfying the following closure and definability axioms.
\begin{enumerate}[label=(\roman*)]
  \item Empty sets are definable:
  \[
    \emptyset\in\Sfam_n.
  \]
  \item Definable sets are closed under finite unions:
  \[
    A,B\in\Sfam_n\quad\Longrightarrow\quad A\cup B\in\Sfam_n.
  \]
  \item Definable sets are closed under finite intersections:
  \[
    A,B\in\Sfam_n\quad\Longrightarrow\quad A\cap B\in\Sfam_n.
  \]
  \item Definable sets are closed under complements:
  \[
    A\in\Sfam_n\quad\Longrightarrow\quad R^n\setminus A\in\Sfam_n.
  \]
  \item Diagonals are definable: for $i<j<n$,
  \[
    \{x\in R^n:x_i=x_j\}\in\Sfam_n.
  \]
  \item Products are definable: if $A\in\Sfam_n$ and $B\in\Sfam_m$, then
  \[
    A\times B
    =\{z\in R^{n+m}:z_L\in A\text{ and }z_R\in B\}
    \in\Sfam_{n+m}.
  \]
  \item One-coordinate projections are definable: if $A\in\Sfam_{n+1}$, then deleting any one
  coordinate gives a set in $\Sfam_n$.
  \item Coordinate reindexing preserves definability: if $A\in\Sfam_n$ and
  $\sigma:\operatorname{Fin}(n)\to\operatorname{Fin}(m)$, then
  \[
    \{x\in R^m:(x_{\sigma(0)},\ldots,x_{\sigma(n-1)})\in A\}
    \in\Sfam_m.
  \]
  \item Block existential projection preserves definability: if $A\in\Sfam_{n+m}$, then
  \[
    \{x\in R^n:\exists y\in R^m,\;x\cat y\in A\}
    \in\Sfam_n.
  \]
  \item The order relation is definable:
  \[
    \{(x,y)\in R^2:x<y\}
    \in\Sfam_2.
  \]
  \item O-minimality in dimension one:
  \[
    A\in\Sfam_1
    \quad\Longleftrightarrow\quad
    A\text{ is a finite union of points and intervals.}
  \]
\end{enumerate}
In Lean this structure is called
\[
  \Lean{OMinimalStructure}.
\]
\end{definition}

\begin{remark}
In many textbook treatments, some of the closure properties above are derived from a smaller list of
axioms.  In the Lean file they are included explicitly.  This is not a mathematical restriction; it is
a practical choice that avoids lengthy low-level proofs about finite coordinates.
\end{remark}

\begin{definition}[Definable set]
Let $M$ be an o-minimal structure.  A subset $A\subseteq R^n$ is called definable if
\[
  A\in\Sfam_n.
\]
In Lean this is abbreviated by
\[
  \Lean{Definable M A}.
\]
\end{definition}

\section{Definable functions}

\subsection{Graphs}

Let
\[
  A\subseteq R^m,
  \qquad
  B\subseteq R^n.
\]
A function $f:A\to B$ is represented in Lean as a function between subtype spaces:
\[
  f:
  \{x:R^m\mid A(x)\}
  \longrightarrow
  \{y:R^n\mid B(y)\}.
\]
Thus an input is not merely a tuple $x\in R^m$, but a pair consisting of $x$ and a proof that
$x\in A$.  Likewise, the output contains both the value $f(x)$ and a proof that this value belongs to
$B$.

The graph of $f$ is a subset of $R^{m+n}$:
\[
  \Graph(f)=\{(x,y)\in R^m\times R^n:x\in A\text{ and }y=f(x)\}.
\]
In terms of concatenated tuples, if $z\in R^{m+n}$, then
\[
  z\in\Graph(f)
  \quad\Longleftrightarrow\quad
  \exists h:A(z_L),\; z_R=f(z_L,h).
\]
This is the Lean definition
\[
  \Lean{FunctionGraph}.
\]

\begin{definition}[Definable function]
A function $f:A\to B$ is definable if
\begin{enumerate}[label=(\roman*)]
  \item $A\in\Sfam_m$;
  \item $B\in\Sfam_n$;
  \item $\Graph(f)\in\Sfam_{m+n}$.
\end{enumerate}
In Lean this bundled object is
\[
  \Lean{DefinableFunction M A B}.
\]
It contains the fields
\[
  \Lean{domain\_mem},\qquad
  \Lean{codomain\_mem},\qquad
  \Lean{toFun},\qquad
  \Lean{graph\_mem}.
\]
\end{definition}

\section{Images of definable functions}

This is the new part added to the Lean file.

\subsection{The image set}

Let
\[
  f:A\to B
\]
be represented as a Lean function
\[
  f:
  \{x:R^m\mid A(x)\}
  \longrightarrow
  \{y:R^n\mid B(y)\}.
\]
The image of $f$ is the subset of $R^n$ defined by
\[
  \Imf(f)=\{y\in R^n:\exists x\in R^m,\exists h:A(x),\;y=f(x,h)\}.
\]
In Lean this is
\[
  \Lean{FunctionImage f}.
\]
Written in ordinary mathematical notation, this is simply
\[
  \Imf(f)=\{y\in R^n:\exists x\in A,\, y=f(x)\}.
\]

\begin{lemma}[The image is contained in the codomain]
For any function $f:A\to B$,
\[
  \Imf(f)\subseteq B.
\]
\end{lemma}

\begin{proof}
Let $y\in\Imf(f)$.  By definition of the image, there exist $x\in R^m$ and a proof $h:A(x)$ such
that
\[
  y=f(x,h).
\]
Since the codomain of $f$ is the subtype $\{y:R^n\mid B(y)\}$, the output $f(x,h)$ comes with a proof
that its underlying tuple belongs to $B$.  Therefore $y\in B$.
\end{proof}

This is the Lean theorem
\[
  \Lean{functionImage\_subset\_codomain}.
\]

\subsection{Why a coordinate swap is needed}

The graph $\Graph(f)$ is stored in coordinates
\[
  (x,y)\in R^{m+n}.
\]
However, the image is obtained by keeping the $y$-coordinates and existentially quantifying over the
$x$-coordinates:
\[
  y\in\Imf(f)
  \quad\Longleftrightarrow\quad
  \exists x\in R^m,
  \;(x,y)
  \in\Graph(f).
\]

The available block projection axiom in the Lean file has the form
\[
  C\in\Sfam_{n+m}
  \quad\Longrightarrow\quad
  \{y\in R^n:\exists x\in R^m,\; y\cat x\in C\}
  \in\Sfam_n.
\]
Thus it projects away the last block.  Since $\Graph(f)$ is in the order $(x,y)$, we first reindex it
into the order $(y,x)$.

Define the block-swap map
\[
  s:R^{n+m}\to R^{m+n}
\]
by
\[
  s(y,x)=(x,y).
\]
In the Lean file this is implemented by two definitions:
\[
  \Lean{imageGraphIndex},
  \qquad
  \Lean{imageGraphArg}.
\]
The theorem
\[
  \Lean{imageGraphArg\_append}
\]
formalises the identity
\[
  s(y\cat x)=x\cat y.
\]

\subsection{Definability of the image}

\begin{theorem}[Definability of images]
Let $M$ be an o-minimal structure.  Let
\[
  A\subseteq R^m,
  \qquad
  B\subseteq R^n,
\]
and let
\[
  f:A\to B
\]
be a definable function.  Then
\[
  \Imf(f)\in\Sfam_n.
\]
Equivalently, the image of a definable function is definable.
\end{theorem}

\begin{proof}
Since $f$ is definable, its graph is definable:
\[
  \Graph(f)
  \in
  \Sfam_{m+n}.
\]

Define the swapped graph
\[
  \Graph(f)^{\mathrm{sw}}
  \subseteq
  R^{n+m}
\]
by
\[
  (y,x)
  \in
  \Graph(f)^{\mathrm{sw}}
  \quad\Longleftrightarrow\quad
  (x,y)
  \in
  \Graph(f).
\]
Equivalently,
\[
  \Graph(f)^{\mathrm{sw}}
  =s^{-1}(\Graph(f)),
\]
where $s(y,x)=(x,y)$ is the block-swap map.

Because coordinate reindexing preserves definability and $\Graph(f)$ is definable, we have
\[
  \Graph(f)^{\mathrm{sw}}
  \in
  \Sfam_{n+m}.
\]

Now apply block existential projection to the last $m$ coordinates.  We obtain the definable set
\[
  P
  =
  \{y\in R^n:\exists x\in R^m,
  \;(y,x)
  \in
  \Graph(f)^{\mathrm{sw}}
  \}
  \in
  \Sfam_n.
\]
By the definition of the swapped graph,
\[
\begin{aligned}
  y\in P
  &\Longleftrightarrow
  \exists x\in R^m,
  \;(y,x)
  \in
  \Graph(f)^{\mathrm{sw}} \\
  &\Longleftrightarrow
  \exists x\in R^m,
  \;(x,y)
  \in
  \Graph(f) \\
  &\Longleftrightarrow
  \exists x\in A,
  \;y=f(x) \\
  &\Longleftrightarrow
  y\in\Imf(f).
\end{aligned}
\]
Therefore $P=\Imf(f)$, and since $P$ is definable, $\Imf(f)$ is definable.
\end{proof}

In Lean the graph-only version is
\[
  \Lean{functionImage\_mem\_of\_graph\_mem}.
\]
The bundled version for a \Lean{DefinableFunction} is
\[
  \Lean{functionImage\_mem}.
\]
The named image of a bundled definable function is
\[
  \Lean{DefinableFunction.image},
\]
and the two final bundled theorems are
\[
  \Lean{DefinableFunction.image\_mem},
  \qquad
  \Lean{DefinableFunction.image\_subset\_codomain}.
\]

\begin{remark}
The assumptions that $A$ and $B$ are definable are included in the bundled notion of
\Lean{DefinableFunction}.  The image-definability proof itself uses only definability of the graph.
The codomain inclusion uses only the fact that the output type is the subtype $B$.
\end{remark}

\section{Composition of definable functions}

\subsection{Relational composition}

Let
\[
  G\subseteq R^{a+b},
  \qquad
  H\subseteq R^{b+c}
\]
be definable relations.  Their relational composition is
\[
  H\circ G
  =
  \{(x,z)\in R^{a+c}:\exists y\in R^b,
  \;(x,y)\in G\text{ and }(y,z)\in H\}.
\]
In Lean this relation is
\[
  \Lean{RelationComp G H}.
\]

\begin{lemma}[Relational composition preserves definability]
If
\[
  G\in\Sfam_{a+b}
  \qquad\text{and}\qquad
  H\in\Sfam_{b+c},
\]
then
\[
  H\circ G\in\Sfam_{a+c}.
\]
\end{lemma}

\begin{proof}
Work in the larger coordinate space
\[
  R^{(a+c)+b}.
\]
We write a point of this space as
\[
  (x,z,y),
\]
where
\[
  x\in R^a,
  \qquad
  z\in R^c,
  \qquad
  y\in R^b.
\]

Using coordinate reindexing, define the two sets
\[
  G'=\bigl\{(x,z,y):(x,y)
  \in G\bigr\},
\]
and
\[
  H'=\bigl\{(x,z,y):(y,z)
  \in H\bigr\}.
\]
Since $G$ and $H$ are definable and reindexing preserves definability, both $G'$ and $H'$ are
definable.  Therefore
\[
  G'\cap H'
\]
is definable.

Project away the final block $y\in R^b$.  The resulting set is
\[
  \{(x,z):\exists y,
  \;(x,y)\in G
  \text{ and }
  (y,z)
  \in H\},
\]
which is exactly $H\circ G$.  Hence $H\circ G$ is definable.
\end{proof}

The Lean proof is
\[
  \Lean{relation\_comp\_mem}.
\]
The auxiliary definitions
\[
  \Lean{compGIndex},\quad
  \Lean{compHIndex},\quad
  \Lean{compGArg},\quad
  \Lean{compHArg}
\]
encode the coordinate reindexings needed to read $(x,y)$ and $(y,z)$ from the larger tuple
$(x,z,y)$.

\subsection{Composition of functions}

\begin{proposition}[Composition of definable functions]
Let
\[
  f:A\to B,
  \qquad
  g:B\to C
\]
be definable functions.  Then
\[
  g\circ f:A\to C
\]
is definable.
\end{proposition}

\begin{proof}
The domain of $g\circ f$ is $A$, which is definable because $f$ is a definable function.  The codomain
is $C$, which is definable because $g$ is a definable function.

It remains to prove that the graph is definable.  For $x\in A$ and $z\in C$,
\[
  (x,z)\in\Graph(g\circ f)
\]
if and only if there exists $y\in B$ such that
\[
  (x,y)\in\Graph(f)
  \qquad\text{and}\qquad
  (y,z)\in\Graph(g).
\]
Thus
\[
  \Graph(g\circ f)
  =
  \Graph(g)\circ\Graph(f).
\]
The right hand side is definable by relational composition.  Hence $g\circ f$ is definable.
\end{proof}

The Lean equality relating relational composition of graphs to the graph of the composed function is
\[
  \Lean{relationComp\_functionGraph\_eq}.
\]
The bundled composition theorem is
\[
  \Lean{DefinableFunction.comp}.
\]

\section{Reusable definability devices for local properties}

The second half of the Lean file proves definability of several one-variable loci.  The proofs all
follow the same pattern:
\begin{enumerate}[label=(\roman*)]
  \item express the desired property using only existential quantifiers, Boolean operations, the graph
  of the function, the order relation, and diagonals;
  \item build a higher-dimensional counterexample set;
  \item project away the witness variables;
  \item take complements to express universal quantifiers;
  \item project again if necessary.
\end{enumerate}

The file introduces several helper predicates.

\subsection{Reindexing a relation into a larger coordinate space}

If $A\subseteq R^n$ and
\[
  \sigma:\operatorname{Fin}(n)\to\operatorname{Fin}(m),
\]
then
\[
  \Lean{HoldsAt A sigma}
\]
is the subset of $R^m$ defined by
\[
  \{z\in R^m:(z_{\sigma(0)},\ldots,z_{\sigma(n-1)})\in A\}.
\]
The theorem
\[
  \Lean{holdsAt\_mem}
\]
says that if $A$ is definable, then \Lean{HoldsAt A sigma} is definable.

\subsection{Domain, graph, order, and equality at chosen coordinates}

For a one-dimensional domain $I\subseteq R$, the predicate
\[
  \Lean{DomainOn I i}
\]
means that the $i$-th coordinate belongs to $I$:
\[
  z_i\in I.
\]

For a graph $G\subseteq R^2$, the predicate
\[
  \Lean{GraphOn G i j}
\]
means
\[
  (z_i,z_j)\in G.
\]

The predicate
\[
  \Lean{LtOn D i j}
\]
means
\[
  z_i<z_j.
\]

The predicate
\[
  \Lean{EqOn i j}
\]
means
\[
  z_i=z_j.
\]
It is definable when $i<j$ by the diagonal axiom, and its complement is used to express inequality.

\subsection{Intervals and cores}

The helper predicate
\[
  \Lean{ValueIntervalAt D a w b}
\]
means
\[
  z_a<z_w<z_b.
\]
The helper predicate
\[
  \Lean{NeighbourhoodAt D c x d}
\]
means
\[
  z_c<z_x<z_d.
\]
Finally,
\[
  \Lean{CoreAt D G x a b v}
\]
means
\[
  (z_x,z_v)\in G
  \qquad\text{and}\qquad
  z_a<z_v<z_b.
\]
Each of these is definable because it is built from reindexed copies of $G$ and of the order relation,
using finite intersections.

The whole space $R^n$ is definable because it is the complement of the empty set.  This is the Lean
theorem
\[
  \Lean{setUniv\_mem}.
\]

\section{Continuity points}

Let
\[
  I\subseteq R
\]
be definable, and let
\[
  f:I\to R
\]
be a definable function.  Write
\[
  G=\Graph(f)\subseteq R^2.
\]

\subsection{The intended mathematical predicate}

A point $x\in R$ is a continuity point of $f$ on $I$ if $x\in I$ and for every open interval
$(a,b)$ containing $f(x)$, there exists an open interval $(c,d)$ containing $x$ such that
\[
  y\in I\text{ and }c<y<d
  \quad\Longrightarrow\quad
  a<f(y)<b.
\]

In graph language this becomes
\[
\begin{aligned}
  x\in I
  \quad\text{and}\quad
  \forall a,b,v\;[&G(x,v)\land a<v<b\Rightarrow{}\\
  &\exists c,d\;(c<x<d\land
  \forall y,w\,(I(y)\land c<y<d\land G(y,w)\Rightarrow a<w<b))].
\end{aligned}
\]
The Lean predicate encoding this mathematical statement is
\[
  \Lean{ContinuousAtOnGraph}.
\]

\subsection{The definable construction}

The Lean proof constructs the discontinuity locus as an existential projection of a definable bad set.
The coordinates in arity $8$ are ordered as
\[
  (x,a,b,v,c,d,y,w).
\]
Define the counterexample set
\[
\begin{aligned}
  C_8=\bigl\{(x,a,b,v,c,d,y,w):{}&
  y\in I,
  \;c<y<d,
  \;G(y,w),\\
  &\neg(a<w<b)
  \bigr\}.
\end{aligned}
\]
This is the Lean definition
\[
  \Lean{Counter8}.
\]
It is definable by the theorem
\[
  \Lean{counter8\_mem}.
\]

Projecting away $(y,w)$ gives a definable subset of $R^6$:
\[
  C_6(x,a,b,v,c,d)
  \quad\Longleftrightarrow\quad
  \exists y,w,
  \;(x,a,b,v,c,d,y,w)
  \in C_8.
\]
This is
\[
  \Lean{CounterExists6}.
\]
Its complement says that the neighbourhood $(c,d)$ has no counterexample.

Next define the good six-dimensional set
\[
\begin{aligned}
  G_6(x,a,b,v,c,d)
  \Longleftrightarrow{}&
  G(x,v)
  \land a<v<b
  \land c<x<d\\
  &\land \neg C_6(x,a,b,v,c,d).
\end{aligned}
\]
This is
\[
  \Lean{Good6}.
\]
Projecting away $(c,d)$ gives
\[
  E_4(x,a,b,v)
  \quad\Longleftrightarrow\quad
  \exists c,d,
  \;G_6(x,a,b,v,c,d).
\]
This is
\[
  \Lean{GoodExists4}.
\]

Now define
\[
\begin{aligned}
  B_4(x,a,b,v)
  \Longleftrightarrow{}&
  G(x,v)
  \land a<v<b
  \land \neg E_4(x,a,b,v).
\end{aligned}
\]
This says that $(a,b)$ is a bad neighbourhood of the value $v=f(x)$.  It is represented by
\[
  \Lean{Bad4}.
\]
Projecting away $(a,b,v)$ gives the bad point set
\[
  B_1(x)
  \quad\Longleftrightarrow\quad
  \exists a,b,v,
  \;B_4(x,a,b,v).
\]
This is
\[
  \Lean{BadPoint1}.
\]
The constructed continuity locus is
\[
  \Cont(f)=I\cap (R\setminus B_1).
\]
In Lean this is
\[
  \Lean{ContinuousPoints}.
\]

\begin{theorem}[Definability of the constructed continuity locus]
If $f:I\to R$ is definable, then \Lean{ContinuousPoints} is definable.
\end{theorem}

\begin{proof}
The set $C_8$ is built from $I$, $G$, the order relation, finite intersections, and a complement.
Hence $C_8$ is definable.  By block projection, $C_6$ is definable.  By complement and finite
intersection, $G_6$ is definable.  By projection, $E_4$ is definable.  By complement and finite
intersection, $B_4$ is definable.  By projection, $B_1$ is definable.  Finally,
\[
  I\cap(R\setminus B_1)
\]
is definable by complement and intersection.
\end{proof}

The Lean theorem is
\[
  \Lean{continuousPoints\_mem}.
\]

\begin{remark}
The Lean file defines both the intended pointwise predicate \Lean{ContinuousAtOnGraph} and the
constructed set \Lean{ContinuousPoints}.  The definability theorem uses the constructed set.  The
mathematical equivalence between the two is exactly the counterexample-unwinding described above.
\end{remark}

\section{Locally constant points}

\subsection{The intended mathematical predicate}

The function $f:I\to R$ is locally constant at $x$ if $x\in I$ and there exists an interval $(c,d)$
containing $x$ such that
\[
  y\in I\text{ and }c<y<d
  \quad\Longrightarrow\quad
  f(y)=f(x).
\]
In graph language this is
\[
\begin{aligned}
  x\in I
  \quad\text{and}\quad
  \exists c,d,v\;[&G(x,v)
  \land c<x<d\\
  &\land
  \forall y,w\,(I(y)\land c<y<d\land G(y,w)
  \Rightarrow w=v)].
\end{aligned}
\]
The corresponding Lean predicate is
\[
  \Lean{LocallyConstantAtOnGraph}.
\]

\subsection{The definable construction}

Use coordinates
\[
  (x,c,d,v,y,w)
  \in R^6.
\]
Define the bad set
\[
\begin{aligned}
  C^{\mathrm{lc}}_6=\bigl\{(x,c,d,v,y,w):{}&
  y\in I,
  \;c<y<d,
  \;G(y,w),
  \;w\ne v
  \bigr\}.
\end{aligned}
\]
This is
\[
  \Lean{LocalConstCounter6}.
\]
It is definable because $w\ne v$ is the complement of the diagonal $w=v$.

Projecting away $(y,w)$ gives
\[
  C^{\mathrm{lc}}_4(x,c,d,v),
\]
which says that the proposed neighbourhood and value admit a counterexample.  This is
\[
  \Lean{LocalConstCounterExists4}.
\]
Its complement says that no such counterexample exists.

Define
\[
\begin{aligned}
  G^{\mathrm{lc}}_4(x,c,d,v)
  \Longleftrightarrow{}&
  G(x,v)
  \land c<x<d
  \land \neg C^{\mathrm{lc}}_4(x,c,d,v).
\end{aligned}
\]
This is
\[
  \Lean{LocalConstGood4}.
\]
Projecting away $(c,d,v)$ gives
\[
  G^{\mathrm{lc}}_1(x)
  \quad\Longleftrightarrow\quad
  \exists c,d,v,
  \;G^{\mathrm{lc}}_4(x,c,d,v).
\]
The final locally constant locus is
\[
  \LC(f)=I\cap G^{\mathrm{lc}}_1.
\]
This is
\[
  \Lean{LocallyConstantPoints}.
\]

\begin{theorem}[Definability of the locally constant locus]
If $f:I\to R$ is definable, then the constructed locally constant locus is definable.
\end{theorem}

\begin{proof}
The set $C^{\mathrm{lc}}_6$ is definable by reindexing $I$, $G$, and the order relation, and by using
the complement of a diagonal to express inequality.  Projection gives $C^{\mathrm{lc}}_4$.  Complement
and intersection give $G^{\mathrm{lc}}_4$.  Projection gives $G^{\mathrm{lc}}_1$.  Intersecting with
$I$ gives $\LC(f)$, which is therefore definable.
\end{proof}

The Lean theorem is
\[
  \Lean{locallyConstantPoints\_mem}.
\]

\section{Locally monotone points}

\subsection{Convention}

The Lean file uses the non-strict convention:
\[
  \text{locally increasing}=\text{locally nondecreasing},
\]
and
\[
  \text{locally decreasing}=\text{locally nonincreasing}.
\]
Since the primitive order symbol is strict, the weak inequality $u\le v$ is expressed as
\[
  \neg(v<u).
\]

\subsection{The intended mathematical predicates}

The function $f:I\to R$ is locally increasing at $x$ if $x\in I$ and there exist $c,d$ such that
$c<x<d$ and for all $y_0,y_1\in I$,
\[
  c<y_0<d,
  \qquad
  c<y_1<d,
  \qquad
  y_0<y_1
  \quad\Longrightarrow\quad
  f(y_0)\le f(y_1).
\]
Equivalently, in graph language, there should be no witnesses $y_0,y_1,w_0,w_1$ such that
\[
  G(y_0,w_0),
  \qquad
  G(y_1,w_1),
  \qquad
  y_0<y_1,
  \qquad
  w_1<w_0.
\]

Similarly, $f$ is locally decreasing at $x$ if there is no such local pair with
\[
  w_0<w_1.
\]
The corresponding Lean predicates are
\[
  \Lean{LocallyIncreasingAtOnGraph},
  \qquad
  \Lean{LocallyDecreasingAtOnGraph}.
\]

\subsection{The common base counterexample set}

Use coordinates
\[
  (x,c,d,y_0,y_1,w_0,w_1)
  \in R^7.
\]
The common base set is
\[
\begin{aligned}
  C^{\mathrm{base}}_7=\bigl\{(x,c,d,y_0,y_1,w_0,w_1):{}&
  y_0\in I,
  \;y_1\in I,
  \;c<y_0<d,
  \;c<y_1<d,\\
  &G(y_0,w_0),
  \;G(y_1,w_1),
  \;y_0<y_1
  \bigr\}.
\end{aligned}
\]
This is
\[
  \Lean{MonotoneBaseCounter7}.
\]
It is definable by finite intersections and reindexing.

The increasing-bad set is
\[
  C^{\mathrm{inc}}_7
  =
  C^{\mathrm{base}}_7
  \cap
  \{w_1<w_0\}.
\]
This is
\[
  \Lean{MonoIncCounter7}.
\]

The decreasing-bad set is
\[
  C^{\mathrm{dec}}_7
  =
  C^{\mathrm{base}}_7
  \cap
  \{w_0<w_1\}.
\]
This is
\[
  \Lean{MonoDecCounter7}.
\]
Both are definable.

\subsection{Projection to good neighbourhoods}

Projecting away
\[
  (y_0,y_1,w_0,w_1)
\]
gives definable subsets of $R^3$:
\[
  C^{\mathrm{inc}}_3(x,c,d),
  \qquad
  C^{\mathrm{dec}}_3(x,c,d).
\]
They say that the interval $(c,d)$ contains a bad pair for increasingness or decreasingness,
respectively.  In Lean these are
\[
  \Lean{MonoIncCounterExists3},
  \qquad
  \Lean{MonoDecCounterExists3}.
\]

Their complements say that the proposed neighbourhood has no bad pair.  Hence define
\[
\begin{aligned}
  G^{\mathrm{inc}}_3(x,c,d)
  &\Longleftrightarrow
  c<x<d
  \land
  \neg C^{\mathrm{inc}}_3(x,c,d),\\
  G^{\mathrm{dec}}_3(x,c,d)
  &\Longleftrightarrow
  c<x<d
  \land
  \neg C^{\mathrm{dec}}_3(x,c,d).
\end{aligned}
\]
These are the Lean definitions
\[
  \Lean{MonoIncGood3},
  \qquad
  \Lean{MonoDecGood3}.
\]

Projecting away $(c,d)$ gives one-dimensional good-point sets
\[
  G^{\mathrm{inc}}_1,
  \qquad
  G^{\mathrm{dec}}_1.
\]
These are
\[
  \Lean{MonoIncGoodPoint1},
  \qquad
  \Lean{MonoDecGoodPoint1}.
\]
The final loci are
\[
  \LInc(f)=I\cap G^{\mathrm{inc}}_1,
  \qquad
  \LDec(f)=I\cap G^{\mathrm{dec}}_1.
\]
In Lean these are
\[
  \Lean{LocallyIncreasingPoints},
  \qquad
  \Lean{LocallyDecreasingPoints}.
\]

\begin{theorem}[Definability of local monotonicity loci]
If $f:I\to R$ is definable, then the constructed locally increasing and locally decreasing loci are
definable.
\end{theorem}

\begin{proof}
The base counterexample set is definable because it is built from $I$, $G$, and $<$ using reindexing
and finite intersections.  Intersecting it with $w_1<w_0$ gives the increasing-bad set, and
intersecting it with $w_0<w_1$ gives the decreasing-bad set.  Both are definable.

Projecting away the variables $y_0,y_1,w_0,w_1$ gives the bad-neighbourhood sets in $R^3$.  Taking
complements expresses that no bad pair exists.  Intersecting with $c<x<d$ and then projecting away
$(c,d)$ gives the points admitting a good neighbourhood.  Finally, intersecting with $I$ gives the
locally increasing and locally decreasing loci.  Each step preserves definability.
\end{proof}

The Lean theorems are
\[
  \Lean{locallyIncreasingPoints\_mem},
  \qquad
  \Lean{locallyDecreasingPoints\_mem}.
\]

\section{Small sanity checks}

The file ends with two simple checks.

\begin{enumerate}[label=(\roman*)]
  \item The empty set is definable in every arity:
  \[
    \Lean{empty\_definable}.
  \]
  \item The order relation is definable:
  \[
    \Lean{order\_definable}.
  \]
\end{enumerate}

\section{Dictionary between Lean names and mathematical objects}

\begin{longtable}{>{\raggedright\arraybackslash}p{0.52\textwidth}>{\raggedright\arraybackslash}p{0.38\textwidth}}
\toprule
Lean name & Mathematical meaning \\
\midrule
\endhead
\Lean{Power R n} & The tuple space $R^n$. \\
\Lean{Power.append} & Tuple concatenation $x\cat y$. \\
\Lean{Power.left}, \Lean{Power.right} & Left and right coordinate blocks of a concatenated tuple. \\
\Lean{setEmpty}, \Lean{setUnion}, \Lean{setInter}, \Lean{setCompl} & Empty set, union, intersection, complement. \\
\Lean{DenseLinearOrderNoEndpoints} & Dense linear order without endpoints. \\
\Lean{Endpoint R} & Formal endpoints $-\infty$, finite points of $R$, and $+\infty$. \\
\Lean{FiniteUnionOfPointsAndIntervals} & Finite unions of points and open intervals in dimension one. \\
\Lean{OMinimalStructure} & The family $(\Sfam_n)_n$ of definable sets with the closure axioms and the one-dimensional o-minimality axiom. \\
\Lean{Definable M A} & The assertion $A\in\Sfam_n$. \\
\Lean{FunctionGraph} & The graph of a function between subtype domains. \\
\Lean{DefinableFunction} & A bundled definable function with definable domain, definable codomain, and definable graph. \\
\Lean{FunctionImage} & The image set $\Imf(f)$. \\
\Lean{imageGraphIndex}, \Lean{imageGraphArg} & The coordinate swap used to turn graph coordinates $(x,y)$ into projection-friendly coordinates $(y,x)$. \\
\Lean{imageGraphArg\_append} & The identity $s(y\cat x)=x\cat y$ for the block-swap map. \\
\Lean{functionImage\_mem\_of\_graph\_mem} & If the graph is definable, then the image is definable. \\
\Lean{functionImage\_mem} & The image of a bundled definable function is definable. \\
\Lean{DefinableFunction.image} & The named image of a bundled definable function. \\
\Lean{DefinableFunction.image\_mem} & The named image is definable. \\
\Lean{DefinableFunction.image\_subset\_codomain} & The image is contained in the codomain. \\
\Lean{RelationComp} & Relational composition $H\circ G$. \\
\Lean{relation\_comp\_mem} & Relational composition preserves definability. \\
\Lean{relationComp\_functionGraph\_eq} & The relational composition of graphs equals the graph of the composed function. \\
\Lean{DefinableFunction.comp} & Composition of definable functions. \\
\Lean{HoldsAt} & Pullback of a definable set by a coordinate map. \\
\Lean{DomainOn} & The condition that a chosen coordinate lies in the domain $I$. \\
\Lean{GraphOn} & The condition that two chosen coordinates form a point of a graph $G$. \\
\Lean{LtOn} & The order condition $z_i<z_j$. \\
\Lean{EqOn} & The equality condition $z_i=z_j$. \\
\Lean{ContinuousPoints} & The constructed continuity locus. \\
\Lean{continuousPoints\_mem} & Definability of the constructed continuity locus. \\
\Lean{LocallyConstantPoints} & The constructed locally constant locus. \\
\Lean{locallyConstantPoints\_mem} & Definability of the constructed locally constant locus. \\
\Lean{LocallyIncreasingPoints} & The constructed locally nondecreasing locus. \\
\Lean{LocallyDecreasingPoints} & The constructed locally nonincreasing locus. \\
\Lean{locallyIncreasingPoints\_mem}, \Lean{locallyDecreasingPoints\_mem} & Definability of the constructed local monotonicity loci. \\
\bottomrule
\end{longtable}

\section{Summary of the new image theorem in one line}

The mathematical content of the new Lean addition is exactly this implication:
\[
  \Graph(f)\in\Sfam_{m+n}
  \quad\Longrightarrow\quad
  \{y\in R^n:\exists x\in R^m,
  \;(x,y)\in\Graph(f)\}
  \in\Sfam_n.
\]
Because
\[
  \{y\in R^n:\exists x\in R^m,
  \;(x,y)\in\Graph(f)\}
  =
  \Imf(f),
\]
this proves that the image of a definable function is definable.

\end{document}
